\section{语音合成}
\label{sec:nlp-speech-tts}

\term{语音合成}（text-to-speech, TTS）把文字及规定的声音条件转换成可播放的波形，要求声音可懂、自然，并且内容正确。同一通知可以快读、慢读、强调新日期或在地点前停顿，这些都是允许的自由；改变日期、数量和否定则不允许。合成面对五项挑战。第一，读法不唯一：“10.5万元”“201”“重庆”在读出之前要先确定发音。第二，长度不匹配：几十个音素要对应几百个声学帧，模型必须决定每个音素持续多久。第三，一对多：同一句文字有许多种合理的语调和节奏，按均方误差回归会把它们平均成一种平淡模糊的声音。第四，相位缺失：梅尔频谱只保留各频带能量，从它恢复波形要补出相位与细节，而且每秒有数万个采样。第五，音色与风格控制：要用指定说话者的声音朗读，最好只给几秒参考录音就能做到。

方法谱系可以按“一个训练好的模型覆盖流水线的哪几段”来看，见图\ref{fig:nlp-speech-tts-lineage}。最早的神经合成系统是两阶段流水线：文本前端处理读法，声学模型把音素映射为梅尔频谱，声码器再把频谱变成波形；Tacotron 2用注意力隐式对齐，FastSpeech 2改为显式预测时长，以解决跳读、重复和串行推断。两阶段分别训练，中间频谱的误差会累积，VITS把声学模型与声码器合成一个条件变分自编码器，联合训练并在内部求对齐。回归式与对抗式目标仍难以表达一对多分布，扩散与流匹配模型从噪声逐步生成频谱，直接学习整个条件分布。最后，神经编解码器把声音变成离散词元，合成被改写为语言模型的续写任务，几秒参考录音就成了上下文提示。节末表\ref{tab:nlp-speech-tts-methods}比较这些方法的对齐方式、自回归性、并行度、零样本能力与内容错误风险。

\begin{figure*}[htbp]
  \centering
  % Four families of TTS systems placed on a common pipeline. Each box spans the
% stages that one trained model covers; arrows mark the interface between models.
\begin{tikzpicture}[
  x=1mm,y=1mm,figure picture,
  stage/.style={font=\figurefont\scriptsize\bfseries,text=FigureInk,align=center},
  fam/.style={font=\figurefont\scriptsize,text=FigureInk,anchor=east,align=right,text width=22mm},
  tm/.style={figure node,minimum height=9mm,inner ysep=2pt,font=\figurefont\scriptsize},
  front/.style={tm,draw=FigureInput,fill=FigureInputFill},
  acou/.style={tm,draw=FigureModel,fill=FigureModelFill},
  mid/.style={tm,draw=FigureStroke,fill=white},
  voc/.style={tm,draw=FigureLoss,fill=FigureLossFill},
  flow/.style={figure flow}
]
  % column headings
  \node[stage] at (14,8) {文本前端};
  \node[stage] at (55,8) {对齐与声学生成};
  \node[stage] at (93,8) {中间表示};
  \node[stage] at (124,8) {波形};
  \draw[figure link,draw=FigureGrid] (0,4) -- (140,4);

  % row 1: two-stage
  \node[fam] at (-2,-6) {两阶段流水线\\（\ref{sec:nlp-speech-acoustics}、\ref{sec:nlp-speech-waveform}）};
  \node[front,text width=24mm] (f1) at (14,-6) {规范化与G2P};
  \node[acou,text width=40mm] (a1) at (55,-6) {Tacotron 2：注意力隐式对齐\\FastSpeech 2：显式时长};
  \node[mid,text width=18mm] (i1) at (93,-6) {梅尔频谱};
  \node[voc,text width=26mm] (v1) at (124,-6) {WaveNet或\\HiFi-GAN};
  \draw[flow] (f1) -- (a1); \draw[flow] (a1) -- (i1); \draw[flow] (i1) -- (v1);

  % row 2: VITS
  \node[fam] at (-2,-20) {端到端VAE\\（\ref{sec:nlp-speech-vits}）};
  \node[front,text width=24mm] (f2) at (14,-20) {音素};
  \node[acou,text width=106mm] (a2) at (85,-20) {VITS：单调对齐搜索与随机时长；规范化流先验；隐变量经对抗训练的生成器直接得到波形};
  \draw[flow] (f2) -- (a2);

  % row 3: diffusion / flow matching
  \node[fam] at (-2,-34) {扩散与流匹配\\（\ref{sec:nlp-speech-diffusion-tts}）};
  \node[front,text width=24mm] (f3) at (14,-34) {音素或字符};
  \node[acou,text width=40mm] (a3) at (55,-34) {Grad-TTS：分数去噪\\F5-TTS：流匹配填充};
  \node[mid,text width=18mm] (i3) at (93,-34) {梅尔频谱};
  \node[voc,text width=26mm] (v3) at (124,-34) {声码器};
  \draw[flow] (f3) -- (a3); \draw[flow] (a3) -- (i3); \draw[flow] (i3) -- (v3);

  % row 4: codec language model
  \node[fam] at (-2,-48) {编解码词元\\语言模型（\ref{sec:nlp-speech-codec-lm-tts}）};
  \node[front,text width=24mm] (f4) at (14,-48) {音素$+$3秒提示};
  \node[acou,text width=40mm] (a4) at (55,-48) {VALL-E：首层自回归\\其余层非自回归};
  \node[mid,text width=18mm] (i4) at (93,-48) {离散声学\\词元};
  \node[voc,text width=26mm] (v4) at (124,-48) {编解码器\\解码器};
  \draw[flow] (f4) -- (a4); \draw[flow] (a4) -- (i4); \draw[flow] (i4) -- (v4);
\end{tikzpicture}

  \caption{四类语音合成系统在同一流水线上的覆盖范围。每个色块是一个单独训练的模型，箭头是模型之间的接口。两阶段流水线和扩散、流匹配模型以梅尔频谱为接口，另配声码器；VITS用一个模型从音素直接生成波形；编解码词元语言模型以离散声学词元为接口，由编解码器的解码器恢复波形。文本前端在端到端系统中可以被字符输入隐式吸收，但读法职责并未消失。}
  \label{fig:nlp-speech-tts-lineage}
\end{figure*}
\FloatBarrier

\subsection{文本前端}
\label{sec:nlp-speech-normalization}

书面文字到发音之间还隔着两步。\term{文本规范化}（text normalization）把日期、金额、编号等非标准写法展开成可以朗读的文字，例如“10.5万元”展开为“十点五万元”；字形到发音映射再决定每个字词的读音，例如“重庆”的“重”读\code{chong2}。直接读取字符的端到端合成没有显式前端，这两项职责由声学模型隐式完成，出错时更难定位，也更难由人工修正。下面依次讨论规则方法、神经规范化与G2P，比较维度是覆盖率、可控性与错误类型。

\subsubsection{规则词典与WFST规范化}
\label{sec:nlp-speech-wfst-normalization}

工业系统的规范化长期采用“先分类、后展开”的两步规则方法。Google的Kestrel系统是一个代表\scite{nlp-speech-ebden2015kestrel}：分类器把输入切成带\term{符号类别}（semiotic class）的单元，如日期、金额、基数、序数、编号、电话、缩写和普通词；展开器再按类别调用相应文法生成读法。两步都写成第\ref{sec:nlp-speech-hybrid}小节识别解码图所用的加权有限状态转导器（WFST）：弧上带输入符号、输出符号和代价，路径把输入串转换为输出串，代价为弧代价之和。

用WFST表达规范化的好处在于可以复合。设$T_{\mathrm{cls}}$把字符串映射为带类别标记的串，$T_{\mathrm{verb}}$把带标记的串映射为读法，两者的复合$T=T_{\mathrm{cls}}\circ T_{\mathrm{verb}}$仍是一个WFST，直接把书面文字映射到读法。对输入串$x$，规范化结果是所有输入为$x$的路径中代价最小者的输出：
\begin{equation}
  \hat y=\operatorname{out}\Bigl(\argmin_{\pi:\,\operatorname{in}(\pi)=x}\;\sum_{e\in\pi}w(e)\Bigr),
  \label{eq:nlp-speech-wfst}
\end{equation}
实现上先构造只接受$x$的线性自动机，与$T$组合后在得到的有向无环图上求最短路径。规则和词典来自显式编写的语言约定，代价可以人工设定，也可以用已标注的书面—读出对估计。

三个例子说明类别怎样决定读法。“10月5日”被日期文法消费为月份10、日期5，输出“十月五日”。“10.5万元”被金额文法输出“十点五万元”，其规范值仍为105000元。“201”有两条候选路径：编号读法“二零一”与基数读法“二百零一”。若前文是“海棠楼”或“房间”，上下文弧给编号路径较低代价，例如编号0.2、基数0.5，最短路径选中“二零一”；若前文是“共有”，代价反过来。同一数字串不能只凭局部形状决定类别。

规则方法可控、可审计：每个读法都能追溯到一条规则，修改一条规则只影响相应类别。缺点是词典缺项与规则冲突需要人工不断补充；规则只看局部窗口，难以利用整句上下文，例如“2/3”在“完成了2/3”中读作分数，在“2/3开会”中读作日期。

\subsubsection{句级神经文本规范化}
\label{sec:nlp-speech-neural-normalization}

规则需要显式枚举写法，整句上下文又会改变某个单元的读法。神经规范化把这件事交给序列到序列模型学习。Zhang等的系统以每个待规范化单元为中心取一个窗口，编码窗口内的字符，由解码器生成该单元的读法；大部分单元原样保留，模型只需输出“不变”标记\scite{nlp-speech-zhang2019textnorm}。Mansfield等改为以整句为输入，使用子词单位，一次生成整句读法，以读取窗口以外的上下文\scite{nlp-speech-mansfield2019}。

训练对$(x,r)$可以来自人工标注，也可以由规则系统生成后筛查；后一种数据会继承规则错误。模型按第\ref{sec:nlp-sft-masked-loss}小节的教师强制最小化$-\sum_j\log p(r_j\mid x,r_{<j})$，推理用贪心或束搜索生成读法。神经模型的典型失败是不可恢复的错误：输出流畅，却把“105000”读成“一万零五百”，听者无从察觉。Zhang等的对策是\term{覆盖文法}（covering grammar）：为每个类别写一个宽松的WFST，只描述“这个数字串可能有哪些读法”，不负责在其中选择；解码时把模型输出限制在覆盖文法接受的集合内，这与第\ref{sec:nlp-inference-grammar}小节的文法约束解码是同一机制。神经模型负责借上下文选择，规则负责保证不越界。

对“预算10.5万元，10月5日确认”，合格的读法应同时保持金额105000元与日期10月5日。可以把输入和输出中的类型、数值与单位分别解析成结构记录，要求两者逐项相等，解析失败的候选不自动接受。这一检查能发现漏读“万”，但发现不了读音错误。正向规范化与第\ref{sec:nlp-speech-itn}小节的反向文本规范化方向相反：“二零一”可能来自“201”也可能来自“二〇一”，两者都不保证可逆。

\subsubsection{字形到发音：G2P}
\label{sec:nlp-speech-g2p}

得到“重庆银行”后，系统还要决定“重”和“行”怎么读。\term{字形到发音}（grapheme-to-phoneme, G2P）把文字映射成发音单位序列。中文先要约定发音单位：可以是带调拼音音节（\code{chong2}），也可以拆为声母、韵母加声调（\code{ch}、\code{ong}、2）；轻声与变调是在前端标注还是交给声学模型实现，也需事先规定。

中文G2P的主要困难在多音字，绝大多数字只有一个读音，查词典即可。对多音字，模型在词典给出的候选集合内做分类。设第$i$个字$x_i$的候选读音集合为$\mathcal{R}(x_i)$，预训练字符编码器给出上下文表示$\bm{h}_i$，则
\begin{equation}
  p(r\mid x,i)=\frac{\exp(\bm{w}_r^{\top}\bm{h}_i)}{\sum_{r'\in\mathcal{R}(x_i)}\exp(\bm{w}_{r'}^{\top}\bm{h}_i)},
  \qquad r\in\mathcal{R}(x_i),
  \label{eq:nlp-speech-polyphone}
\end{equation}
分母只在候选集合上求和，模型不会给“行”输出\code{chong2}这类不可能的读音。g2pM用双向LSTM实现这一分类，并给出了一个人工标注的普通话多音字基准\scite{nlp-speech-park2020g2pm}。“重庆”为\code{chong2 qing4}，“银行”为\code{yin2 hang2}；把每个字最常见的读音当成唯一读音，就会读出“zhong4庆”“银xing2”。

词典之外的词，尤其是外文名和新词，需要生成而非分类。Yao与Zweig把G2P看作字母序列到音素序列的翻译，用编码器—解码器模型以交叉熵训练\scite{nlp-speech-yao2015g2p}，推理在音素词表上束搜索，再核对音节数与语种标识。拿到本调读音后，再按规则处理变调，例如两个上声相连时前一个读作阳平：“展览”本调为\code{zhan3 lan3}，实际读作\code{zhan2 lan3}。

通知里的“海棠楼”核为\code{hai3 tang2 lou2}，“201”先经规范化得到“二零一”再转为发音单位；\code{ML-0418}若要朗读，必须事先规定字母和数字的逐位读法。G2P只决定读哪个音，停顿、句重音和语速由后面的声学模型决定。

\subsection{两阶段声学模型}
\label{sec:nlp-speech-acoustics}

两阶段流水线中的声学模型把发音序列$x=(x_1,\ldots,x_N)$映射为梅尔频谱$\bm{M}\in\mathbb{R}^{T\times80}$，每行是一帧80维的对数梅尔能量（第\ref{sec:nlp-speech-features}小节）；声码器再从$\bm{M}$生成波形。声学模型的核心困难是长度差：$N$是几十，$T$是几百，而且每个音素对应多少帧事先未知。下面两个模型给出两种回答。Tacotron 2逐帧自回归生成，用注意力在生成过程中隐式决定当前帧读哪个字符；FastSpeech 2先显式预测每个音素的帧数，展开后并行生成全部帧。频谱提取参数（采样率、帧移、梅尔范围）须与声码器一致，不能只因两边都是80个频带就直接拼接。

\subsubsection{注意力声文对齐：Tacotron 2}
\label{sec:nlp-speech-tacotron}

Tacotron 2把合成看作序列到序列生成：编码器读字符，解码器每步生成一帧梅尔频谱，注意力决定这一帧对应哪些字符\scite{nlp-intro-shen2018}。图\ref{fig:nlp-speech-tacotron2}给出完整结构，下面按数据流逐层说明。

\begin{figure*}[tp]
  \centering
  % Tacotron 2 architecture: encoder (left), location-sensitive attention (middle),
% autoregressive decoder with post-net (right). Layer sizes follow Shen et al. (2018).
\begin{tikzpicture}[
  x=1mm,y=1mm,figure picture,
  enc/.style={figure node,draw=FigureInput,fill=FigureInputFill,text width=33mm,minimum height=9mm},
  dec/.style={figure node,draw=FigureModel,fill=FigureModelFill,text width=33mm,minimum height=9mm},
  att/.style={figure node,draw=FigureLoss,fill=FigureLossFill,text width=27mm,minimum height=11mm},
  io/.style={figure node,draw=FigureStroke,fill=white,text width=24mm,minimum height=8mm},
  flow/.style={figure flow},
  tstep/.style={figure flow,dash pattern=on 2pt off 1.5pt,draw=FigureMuted},
  lab/.style={figure edge label,font=\figurefont\scriptsize}
]
  % encoder
  \node[io] (chars) at (0,0) {字符序列，长度$N$};
  \node[enc] (emb) at (0,14) {字符嵌入\\$N\times512$};
  \node[enc] (conv) at (0,28) {3层卷积\\核宽5，512通道};
  \node[enc] (lstm) at (0,42) {双向LSTM\\每向256单元};
  \node[io,text width=30mm] (H) at (0,56) {编码$\bm H\in\mathbb R^{N\times512}$};
  \foreach \a/\b in {chars/emb,emb/conv,conv/lstm,lstm/H}{\draw[flow] (\a) -- (\b);}

  % attention
  \node[att] (att) at (50,42) {位置敏感注意力\\权重$\bm\alpha_i$};
  \draw[flow] (H.east) -- (att.north west);
  \draw[tstep] ([xshift=4mm]att.north) .. controls +(1,9) and +(9,7) .. (att.north east);
  \node[anchor=south,font=\figurefont\scriptsize,text=FigureMuted,align=center] at (57,55.5) {累积权重$\sum_{k<i}\bm\alpha_k$\\经卷积成位置特征};

  % decoder
  \node[io,text width=30mm] (prev) at (100,0) {前一帧$\bm m_{i-1}\in\mathbb R^{80}$};
  \node[dec] (pre) at (100,14) {预网络\\2层全连接256};
  \node[dec] (dlstm) at (100,28) {2层单向LSTM\\每层1024单元};
  \node[dec] (proj) at (100,42) {两个线性投影\\输入$[\bm s_i;\bm c_i]$};
  \draw[flow] (prev) -- (pre);
  \draw[flow] (pre) -- (dlstm);
  \draw[flow] (dlstm) -- (proj);
  \draw[flow] ([yshift=1.5mm]dlstm.west) -- (60,29.5) -- node[lab,pos=.45] {$\bm s_i$} (60,36.5);
  \draw[flow] (att.east) -- node[lab] {$\bm c_i$} (proj.west);
  \draw[tstep] (44,36.5) |- node[lab,pos=.75] {$\bm c_{i-1}$} ([yshift=-2mm]dlstm.west);

  % outputs
  \node[io,text width=20mm] (stop) at (80,62) {停止概率$\hat p_i$};
  \node[io,text width=20mm] (mel) at (114,62) {梅尔帧$\hat{\bm m}_i$};
  \draw[flow] (proj.north) ++(-6,0) |- ++(0,4) -| (stop.south);
  \draw[flow] (proj.north) ++(6,0) |- ++(0,4) -| (mel.south);
  \node[dec,text width=28mm] (post) at (114,74) {后处理网络\\5层卷积，预测残差};
  \node[circle,draw=FigureStroke,line width=.75pt,inner sep=1pt] (plus) at (114,84) {$+$};
  \node[io,text width=30mm] (melp) at (114,93) {修正频谱$\widehat{\bm M}^{+}$};
  \draw[flow] (mel) -- (post);
  \draw[flow] (post) -- (plus);
  \draw[flow] (plus) -- (melp);
  \draw[flow] (mel.west) -- (95,62) |- (plus.west);
  \draw[tstep] (mel.east) -- (134,62) |- node[pos=.25,fill=white,font=\figurefont\scriptsize,text=FigureMuted,align=center] {推断时\\反馈；\\训练时\\用真实帧} (prev.east);
\end{tikzpicture}

  \caption{Tacotron 2的结构。左列为编码器：长度为$N$的字符序列经嵌入、3层卷积和双向LSTM得到编码$\bm H$。中间为位置敏感注意力，它在打分时读取此前各步注意力权重的累积和。右列为解码器：前一帧经预网络，与上一步上下文$\bm c_{i-1}$一起送入2层LSTM，得到状态$\bm s_i$；$[\bm s_i;\bm c_i]$经两个线性投影分别输出梅尔帧与停止概率，后处理网络预测残差得到修正频谱。虚线表示跨时间步的连接。层数与宽度取自原论文。}
  \label{fig:nlp-speech-tacotron2}
\end{figure*}

编码器把字符序列变成上下文相关的表示。每个字符先查512维嵌入，得到$N\times512$的矩阵；接着经过3层一维卷积，每层512个核宽为5的滤波器，后接批归一化与ReLU，卷积让每个位置看到前后各若干字符，起到类似$n$元特征的作用；最后一层双向LSTM每个方向256单元，拼接后得到$\bm{H}=(\bm{h}_1,\ldots,\bm{h}_N)\in\mathbb{R}^{N\times512}$。

解码器第$i$步要预测第$i$帧$\bm{m}_i\in\mathbb{R}^{80}$。前一帧$\bm{m}_{i-1}$先经\term{预网络}（pre-net），即两层256单元的全连接层加ReLU，并以0.5的概率丢弃；丢弃在推断时也保持开启，原论文用它给输出引入随机性。预网络输出与上一步的上下文向量$\bm{c}_{i-1}$拼接，送入两层各1024单元的单向LSTM，得到解码器状态$\bm{s}_i$。

注意力用$\bm{s}_i$在$\bm{H}$上取上下文。普通加性注意力只比较$\bm{s}_i$与$\bm{h}_j$的内容，长句中容易跳到内容相似的远处字符。\term{位置敏感注意力}（location-sensitive attention）把“此前已经读到哪里”作为额外特征\scite{nlp-speech-chorowski2015}。设此前各步注意力权重的累积和为$\bm{\beta}_{i}=\sum_{k<i}\bm{\alpha}_k\in\mathbb{R}^{N}$，用$32$个长度为31的一维卷积核$\bm{F}$与它卷积，得到每个位置$j$的位置特征$\bm{f}_{i,j}\in\mathbb{R}^{32}$，则
\begin{equation}
\begin{aligned}
  e_{i,j}&=\bm{v}^{\top}\tanh\bigl(\bm{W}\bm{s}_i+\bm{V}\bm{h}_j+\bm{U}\bm{f}_{i,j}+\bm{b}\bigr),\qquad
  \bm{f}_{i}=\bm{F}*\bm{\beta}_{i},\\
  \alpha_{i,j}&=\frac{\exp(e_{i,j})}{\sum_{j'=1}^{N}\exp(e_{i,j'})},\qquad
  \bm{c}_i=\sum_{j=1}^{N}\alpha_{i,j}\bm{h}_j .
\end{aligned}
\label{eq:nlp-speech-location-attention}
\end{equation}
$\bm{W}$、$\bm{V}$、$\bm{U}$把三种输入投影到128维，$\bm{v}\in\mathbb{R}^{128}$给出标量分数。累积权重记录了每个字符已被读过多少：已读完的字符累积值接近1，尚未读到的接近0，卷积核可以学会“偏好累积值刚从1降到0的位置”，即紧接着已读部分的下一个字符。这样，对齐倾向于单调向前移动。

输出端把$[\bm{s}_i;\bm{c}_i]$经一个线性投影得到梅尔帧$\hat{\bm{m}}_i$，经另一个线性投影和sigmoid得到\term{停止标记}（stop token）概率$\hat p_i$，表示合成是否应在这一帧结束。全部帧拼成$\widehat{\bm{M}}$后，\term{后处理网络}（post-net）用5层一维卷积（每层512个核宽为5的滤波器，批归一化，除最后一层外接tanh）预测残差，得到$\widehat{\bm{M}}^{+}=\widehat{\bm{M}}+\operatorname{PostNet}(\widehat{\bm{M}})$。后处理网络能看到整段频谱，可以修正逐帧生成时只看过去造成的局部误差。训练目标为
\begin{equation}
  L_{\mathrm{Taco}}=\lVert\widehat{\bm{M}}-\bm{M}\rVert_2^2
  +\lVert\widehat{\bm{M}}^{+}-\bm{M}\rVert_2^2
  +\sum_{i=1}^{T}\operatorname{BCE}(\hat p_i,\,\mathbb{1}[i=T]),
  \label{eq:nlp-speech-tacotron-loss}
\end{equation}
两项频谱误差分别监督后处理前后的输出，第三项是停止标记的二元交叉熵，只有最后一帧的目标为1。训练时解码器输入真实的前一帧（教师强制），推断时输入自己生成的帧；推断在$\hat p_i$首次超过0.5时停止，实际系统还应设最大帧数，并把截断标记为异常而非正常完成。原系统的声码器是以预测频谱为条件的修改版WaveNet，见第\ref{sec:nlp-speech-wavenet}小节。

Tacotron 2的优点是结构统一，对齐不需要外部标注，音质在当时接近真人录音。它的缺点都来自隐式对齐与自回归。注意力可能卡在某个字符附近导致重复，也可能提前跳到句尾导致漏读，长句与训练集中少见的文本尤其容易出错。逐帧生成使推断时间与帧数成正比，无法并行。训练用真实帧、推断用自己生成的帧，误差会沿时间累积。对“四月十八日”，可以观察注意力是否依次经过各字，但注意力图只能提示问题，是否漏读仍需听觉核验。

\subsubsection{时长与韵律条件：FastSpeech 2}
\label{sec:nlp-speech-fastspeech}

Tacotron 2的跳读、重复与串行推断都源于“长度在生成中隐式形成”。若先给每个音素分配帧数，就能在长度确定后一次性并行生成全部帧，对齐也不会跳读。FastSpeech用自回归教师模型的注意力提取时长，并以教师合成的频谱作为蒸馏目标\scite{nlp-speech-ren2019fastspeech}；FastSpeech 2去掉教师，改用强制对齐器从真实录音提取时长，直接以真实频谱为目标，并把音高与能量也作为显式条件输入，以缓解一对多问题\scite{nlp-speech-ren2021}。

模型由三部分组成：音素编码器、方差适配器和梅尔频谱解码器。编码器与解码器各由4个\term{前馈Transformer块}（feed-forward Transformer block, FFT块）堆叠而成。FFT块与第\ref{sec:nlp-llm-decoder-layer}小节的Transformer层相同，只是把逐位置前馈网络换成两层一维卷积，以便利用相邻帧或相邻音素的局部关系。设输入$\bm{X}\in\mathbb{R}^{L\times256}$，
\begin{equation}
\begin{aligned}
  \bm{X}'&=\operatorname{LN}\bigl(\bm{X}+\operatorname{MHSA}(\bm{X})\bigr),\\
  \operatorname{FFT}(\bm{X})&=\operatorname{LN}\bigl(\bm{X}'+\operatorname{Conv}_{2}(\operatorname{ReLU}(\operatorname{Conv}_{1}(\bm{X}')))\bigr),
\end{aligned}
\label{eq:nlp-speech-fft-block}
\end{equation}
$\operatorname{MHSA}$为多头自注意力，$\operatorname{LN}$为层归一化；$\operatorname{Conv}_1$把256通道升到1024通道，$\operatorname{Conv}_2$再降回256通道。原文配置中自注意力为2头，两层卷积的核宽分别为9和1。编码器输入是音素嵌入加正弦位置编码，长度$L=N$；解码器输入是方差适配器展开后的序列，长度$L=T$，最后经线性层输出$T\times80$的频谱。

\term{方差适配器}（variance adaptor）依次加入时长、音高和能量三类信息。三个预测器结构相同：两层一维卷积（核宽3、256通道、ReLU），每层后接层归一化与丢弃，最后经线性层为每个位置输出一个标量。时长预测器作用在音素级编码$\bm{h}_1,\ldots,\bm{h}_N$上，预测对数时长$\hat r_i\approx\log(1+d_i)$，其中$d_i$是第$i$个音素的整数帧数；训练目标$d_i$由Montreal强制对齐器把音素对到录音后换算得到\cite{nlp-speech-mcauliffe2017}，强制对齐的原理见第\ref{sec:nlp-speech-forced-alignment}小节。\term{长度调节器}按$d_i$重复第$i$个向量，使隐藏长度与频谱一致：
\begin{equation}
 \widetilde{\bm{h}}_t=\bm{h}_i
 \quad\text{当}\quad
 \sum_{j<i}d_j<t\leq\sum_{j\leq i}d_j,
 \qquad T=\sum_{i=1}^{N}d_i.
 \label{eq:nlp-speech-length-regulator}
\end{equation}
教学音素序列/p, a, t, a/的时长为$(2,3,1,2)$帧，展开为
\[
  (\bm{h}_1,\bm{h}_1,\bm{h}_2,\bm{h}_2,\bm{h}_2,\bm{h}_3,\bm{h}_4,\bm{h}_4),
\]
共8帧；帧移10\,ms时对应80\,ms。推断时用$\hat d_i=\max(1,\operatorname{round}(\exp\hat r_i-1))$恢复整数，下界1保证非静音音素不被删掉，这是实现约定。语速控制把时长除以速度因子再取整：对$(2,3,1,2)$取因子1.5，得$(1.33,2,0.67,1.33)$，四舍五入并取下界1后为$(1,2,1,1)$，共5帧；$8/1.5$不是整数，实际加速比也不恰好是1.5。

长度调节之后，音高与能量预测器在帧级序列$\widetilde{\bm{H}}\in\mathbb{R}^{T\times256}$上工作。音高从训练录音估计基频得到，无声帧用插值补齐；能量取每帧STFT幅度向量的$L_2$范数。原方法还把连续音高曲线作连续小波变换，预测其多尺度表示后再逆变换恢复，以便刻画不同时间尺度的起伏。为把标量条件加回隐藏序列，FastSpeech 2把它们分桶量化后查嵌入表：取边界$b_0<b_1<\cdots<b_{K}$（音高在对数刻度上均匀取，能量在线性刻度上均匀取，$K=256$），第$t$帧的值$p_t$落入区间$[b_{k-1},b_k)$时取第$k$个嵌入$\bm{E}^{\mathrm{pitch}}_k\in\mathbb{R}^{256}$，
\begin{equation}
  \widetilde{\bm{h}}_t\leftarrow\widetilde{\bm{h}}_t
  +\bm{E}^{\mathrm{pitch}}_{\operatorname{bin}(p_t)}
  +\bm{E}^{\mathrm{energy}}_{\operatorname{bin}(\varepsilon_t)} .
  \label{eq:nlp-speech-variance-embed}
\end{equation}
以教学化的4个能量桶为例，边界取0、25、50、75、100，能量值12、48、51、80分别落入第1、2、3、4桶；48与51数值接近却落入不同桶，这是量化的代价，256个桶时相邻桶的嵌入在训练中会变得相近。训练时方差适配器使用从录音提取的真实时长、音高与能量（与教师强制同理），推断时全部改用预测值。

总目标是频谱的$L_1$误差加三项方差预测的均方误差：
\begin{equation}
  L_{\mathrm{FS2}}=\lVert\widehat{\bm{M}}-\bm{M}\rVert_1
  +\sum_{i}\bigl(\hat r_i-\log(1+d_i)\bigr)^2
  +\sum_{t}(\hat p_t-p_t)^2+\sum_{t}(\hat\varepsilon_t-\varepsilon_t)^2 .
  \label{eq:nlp-speech-fs2-loss}
\end{equation}
显式的音高与能量条件解决了一部分一对多问题：给定这些值后，频谱的不确定性大大降低，解码器不必再把不同语调平均成一种。但三个预测器本身仍以均方误差回归，预测出的音高曲线依然趋于平均，这是后面扩散与流匹配模型要解决的问题。

与Tacotron 2相比，FastSpeech 2的对齐由时长决定，不会跳读或重复，全部帧并行生成，语速、音高和能量可以直接调节，例如对通知强调“十八日”时延长相应音素并提高音高，而不改变“校外来宾”和“前”等内容。代价是依赖外部强制对齐器提供时长，对齐器的错误会直接成为训练目标；时长取整和分桶量化也引入了离散误差。

\subsection{神经声码器}
\label{sec:nlp-speech-waveform}

梅尔频谱只保留各频带的能量，丢掉了相位和频带内的细节；\term{声码器}（vocoder）要在给定频谱的条件下补出这些信息，生成波形。一个可复现的接口至少规定频谱提取参数、采样率、每帧对应的采样数（帧移）以及补边与裁剪方式；波形时长由采样数除以采样率得到，不能把帧数直接当成秒数。设采样率24\,kHz、帧移12.5\,ms，则每帧对应300个采样，80帧对应24000个采样，即1秒。声码器的谱系从逐采样自回归的WaveNet开始；为了并行生成，后续出现了知识蒸馏的并行WaveNet、基于规范化流的WaveGlow，以及基于生成对抗网络的MelGAN与HiFi-GAN。下面写出两端的代表：WaveNet与HiFi-GAN。

\subsubsection{条件逐采样生成：WaveNet}
\label{sec:nlp-speech-wavenet}

WaveNet把波形看成一维序列，按自回归分解其概率\scite{nlp-speech-oord2016wavenet}。设目标波形为$x_1,\ldots,x_N$，梅尔条件为$\bm{M}$，则
\begin{equation}
 p(x_{1:N}\mid\bm{M})
   =\prod_{n=1}^{N}p_\theta(x_n\mid x_{<n},\bm{M}),
 \qquad
 L_{\mathrm{wave}}=-\sum_{n=1}^{N}
       \log p_\theta(x_n\mid x_{<n},\bm{M}).
 \label{eq:nlp-speech-wavenet}
\end{equation}
困难在于$x_n$依赖很长的历史：一个音节持续数千个采样。WaveNet用\term{空洞因果卷积}（dilated causal convolution）扩大感受野。因果指输出只依赖当前及过去的采样；空洞率为$r$、核宽为$K$的卷积把相隔$r$的采样组合起来，$y_n=\sum_{k=0}^{K-1}w_k\,x_{n-kr}$。逐层把空洞率加倍，感受野呈指数增长。叠$L$层、第$\ell$层空洞率为$r_\ell$时，感受野为
\begin{equation}
  R=1+(K-1)\sum_{\ell=1}^{L}r_\ell .
  \label{eq:nlp-speech-receptive-field}
\end{equation}
取$K=2$，空洞率$1,2,4,\ldots,512$共10层为一组，一组的感受野为$1+1023=1024$个采样；重复3组得$1+3\times1023=3070$个采样，24\,kHz下约128毫秒。若不用空洞而要达到同样感受野，需要三千多层普通因果卷积。

每层的非线性是\term{门控激活单元}（gated activation unit）。设第$\ell$层输入为$\bm{z}$，上采样到采样率的梅尔条件为$\bm{y}$，则
\begin{equation}
  \bm{o}=\tanh(\bm{W}_{f}*\bm{z}+\bm{V}_{f}\bm{y})\odot\sigma(\bm{W}_{g}*\bm{z}+\bm{V}_{g}\bm{y}),
  \label{eq:nlp-speech-wavenet-gate}
\end{equation}
$*$是空洞因果卷积，$\bm{V}_f\bm{y}$、$\bm{V}_g\bm{y}$是条件的$1\times1$卷积，$\odot$为逐元素乘。tanh支路产生候选信号，sigmoid支路决定放行多少，与LSTM的门同理。$\bm{o}$经$1\times1$卷积后一路加回输入形成残差连接，一路作为跳跃连接汇总到输出端；所有层的跳跃输出相加，经两次ReLU与$1\times1$卷积后得到输出分布。梅尔条件$\bm{y}$由转置卷积把每帧上采样300倍，与采样对齐。

输出分布有两种。原WaveNet把16位采样先用\term[mu律压扩]{μ律压扩}（μ-law companding）压缩，再均匀量化为256类，用softmax建模：
\begin{equation}
  f(x)=\operatorname{sign}(x)\frac{\ln(1+\mu\lvert x\rvert)}{\ln(1+\mu)},\qquad
  x\in[-1,1],\ \mu=255 .
  \label{eq:nlp-speech-mulaw}
\end{equation}
它给小振幅更细的分辨率：$x=0.01$映射为$\ln3.55/\ln256\approx0.228$，$x=0.5$映射为$\ln128.5/\ln256\approx0.876$；在正半轴上，原始区间$[0,0.01]$只占1\%，压扩后却占约23\%的量化级。人耳对安静段的量化噪声更敏感，压扩把量化级更多地分给它们。Tacotron 2中的修改版改为输出10个分量的混合逻辑分布参数，直接对连续采样值计算似然，并以Tacotron 2在教师强制下预测的频谱为条件训练，使训练条件接近推断时的上游输出。

WaveNet的优点是音质高、概率模型明确；缺点是推断必须逐采样串行，生成1秒24\,kHz声音需要24000步，每步都要过一遍全部层，缓存可以避免重复计算，却不能把串行变成并行。另一方面，声码器只负责把频谱变成声音：上游频谱若对应“十七日”，再自然的波形也不会变成“十八日”。

\subsubsection{并行波形生成：HiFi-GAN}
\label{sec:nlp-speech-hifigan}

要并行生成波形，可以用一个卷积网络直接把整段梅尔频谱上采样到波形长度。难点在训练目标：只用频谱或波形的$L_1$误差训练，会得到相位模糊、带金属声的结果，因为满足同一频谱的波形有许多个，回归目标会把它们平均。第\ref{sec:nlp-speech-soundstream}小节的编解码器已经用对抗训练解决同一问题：判别器学习区分真实与生成的样本，生成器学习让判别器无法区分，于是生成样本被推向真实分布的某一个模式而非平均。HiFi-GAN把这一思路落实为一个生成器和两组针对语音周期结构设计的判别器\scite{nlp-speech-kong2020}。

生成器输入$\bm{M}\in\mathbb{R}^{T\times80}$，先经一层核宽7的卷积升到512通道，再经4次转置卷积上采样。以22.05\,kHz、帧移256个采样的配置为例，四次上采样倍率为$(8,8,2,2)$，乘积为256，通道依次减半为256、128、64、32，长度从$T$变为$256T$。每次上采样之后接一个\term{多感受野融合}（multi-receptive field fusion, MRF）模块：它并行使用核宽为3、7、11的三个残差块，每个残差块内部用空洞率1、3、5的空洞卷积，输出取三者的平均。不同核宽与空洞率覆盖不同长度的模式，MRF让每一层同时看到短周期与长周期结构。最后一层卷积接tanh，输出$[-1,1]$内的波形。

\term{多周期判别器}（multi-period discriminator, MPD）针对语音的周期性。浊音由基频决定的周期重复，同一相位上的采样之间有规律。MPD把长度为$N$的波形按周期$p$重排为$\lceil N/p\rceil\times p$的二维矩阵，第$a$行第$b$列为$x_{(a-1)p+b}$，再用沿第一维的二维卷积（核宽$(5,1)$）处理，每一列只包含间隔恰为$p$的采样。例如$x_1,\ldots,x_{12}$取$p=3$，三列分别为$(x_1,x_4,x_7,x_{10})$、$(x_2,x_5,x_8,x_{11})$、$(x_3,x_6,x_9,x_{12})$。HiFi-GAN取$p\in\{2,3,5,7,11\}$五个子判别器，用互质的周期尽量避免重叠。\term{多尺度判别器}（multi-scale discriminator, MSD）沿用MelGAN的设计\cite{nlp-speech-kumar2019melgan}，在原波形以及2倍、4倍平均池化后的波形上各用一个一维卷积判别器，检查连续的局部模式与较长的时间结构；MPD只看等间隔采样，MSD补上了相邻采样之间的关系。

训练采用最小二乘形式的对抗损失。令$\hat x=G(\bm{M})$，$D_k$为第$k$个子判别器（MPD与MSD共8个），$\phi$为固定的梅尔提取器，$D_k^{(l)}$为$D_k$第$l$层的特征、共$N_l$个元素：
\begin{equation}
\begin{aligned}
 L_D&=\sum_k\mathbb{E}\bigl[(D_k(x)-1)^2+D_k(\hat x)^2\bigr],\\
 L_G&=\sum_k\Bigl[\mathbb{E}(D_k(\hat x)-1)^2+\lambda_{\mathrm{fm}}L_{\mathrm{fm}}(G;D_k)\Bigr]\\
 &\quad+\lambda_{\mathrm{mel}}\,\mathbb{E}\lVert\phi(x)-\phi(\hat x)\rVert_1,\\
 L_{\mathrm{fm}}(G;D_k)&=\mathbb{E}\sum_{l}\frac{1}{N_l}\bigl\lVert D_k^{(l)}(x)-D_k^{(l)}(\hat x)\bigr\rVert_1,
\end{aligned}
\label{eq:nlp-speech-hifigan}
\end{equation}
系数取$\lambda_{\mathrm{fm}}=2$，$\lambda_{\mathrm{mel}}=45$。判别器希望真实样本输出1、生成样本输出0；生成器希望生成样本输出1。特征匹配损失与EnCodec中的同名项作用相同（第\ref{sec:nlp-speech-soundstream}小节），比较真实与生成波形在判别器各中间层的特征，提供比单个真假分数更稠密的梯度；梅尔损失约束生成波形的频谱与条件一致，保证内容不走样，并使训练初期稳定。两者交替更新，推断只保留生成器，一次前向即得到整段波形。

HiFi-GAN的优点是并行生成，在GPU上远快于实时，音质与自回归声码器相当。代价是训练依赖判别器设计与损失系数，对抗训练可能不稳定，生成波形偶尔出现周期性伪影；判别器认为逼真，也不说明每个数字都读对了。比较WaveNet与HiFi-GAN的速度时，应统一采样率、频谱定义、硬件、批量和计时边界，分别记录首段可播放时间与整段生成时间。

沿前端、声学模型与声码器贯通一份通知：“机器学习读书会改至4月18日14:00，地点为海棠楼201。校外来宾须在4月16日18:00前提交报名信息。”规范化保留活动时间与报名截止两个字段，产生“四月十八日下午两点”与“四月十六日十八点前”等读法；G2P核对“海棠楼”及编号；声学模型保留每个音节和断句；声码器按约定采样率恢复波形。最终声音按四项独立问题验收：是否听到4月18日14:00，是否听到海棠楼201，是否保留“校外来宾”这一对象，是否听到严格早于4月16日18:00的截止条件。回转识别可以自动比较这些字段，但识别器可能把发错的音纠正成常用词，因此只作筛查；真人盲听转写、配对偏好和平均意见分（mean opinion score, MOS）另行记录材料、评价者与量表，内容正确、可懂度和自然度分别报告。

\subsection{端到端合成：VITS}
\label{sec:nlp-speech-vits}

两阶段流水线有两个遗留问题。声学模型与声码器分别训练，声码器在训练时看到的是真实频谱，推断时却要处理声学模型的预测频谱，两者的误差累积；FastSpeech 2又依赖外部对齐器提供时长。VITS用一个模型同时解决这两个问题：它把从文字到波形的生成写成条件变分自编码器，以隐变量替代梅尔频谱作为内部接口，对齐由模型在训练中自己求得\scite{nlp-speech-kim2021vits}。

先回顾所需的两个工具。变分自编码器（VAE）为观测$\bm{x}$引入隐变量$\bm{z}$，用后验近似$q_\phi(\bm{z}\mid\bm{x})$与先验$p_\theta(\bm{z})$构造对数似然的下界：$\log p_\theta(\bm{x})\geq\mathbb{E}_{q_\phi}[\log p_\theta(\bm{x}\mid\bm{z})]-\operatorname{KL}(q_\phi(\bm{z}\mid\bm{x})\,\Vert\,p_\theta(\bm{z}))$，即\term{证据下界}，第一项要求能从$\bm{z}$重建$\bm{x}$，第二项要求后验贴近先验。规范化流用一个可逆映射$f$把简单分布变成复杂分布，换元后密度为$p(\bm{z})=\mathcal{N}(f(\bm{z}))\,\lvert\det\partial f(\bm{z})/\partial\bm{z}\rvert$，可精确计算。两者以及后两小节要用的扩散模型与流匹配，其一般推导可参考\BookChapterRef{chap:pgm-generative}{卷六的概率生成模型章}；本小节只用上述两个结论。

VITS的条件VAE由四部分组成。后验编码器读取目标语音的线性频谱$\bm{x}_{\mathrm{lin}}$（比梅尔频谱保留更多细节），经16个非因果WaveNet残差块输出每帧的均值与方差，给出$q_\phi(\bm{z}\mid\bm{x}_{\mathrm{lin}})$，$\bm{z}\in\mathbb{R}^{T\times192}$。解码器就是HiFi-GAN的生成器，把$\bm{z}$直接上采样为波形；为节省显存，训练时只解码随机截取的一小段$\bm{z}$。先验编码器由6层Transformer文本编码器和一个规范化流$f_\theta$组成：文本编码器为每个音素$i$输出均值$\bm{\mu}_i$与标准差$\bm{\sigma}_i$，对齐矩阵$A$把它们分配到各帧；流由4个仿射耦合层构成，把后验样本变换到文本给出的高斯分布上：
\begin{equation}
  p_\theta(\bm{z}\mid c,A)=\mathcal{N}\bigl(f_\theta(\bm{z});\,\bm{\mu}_{\theta}(c,A),\,\bm{\sigma}_{\theta}(c,A)\bigr)\,\Bigl\lvert\det\frac{\partial f_\theta(\bm{z})}{\partial\bm{z}}\Bigr\rvert,
  \label{eq:nlp-speech-vits-prior}
\end{equation}
$c$是音素序列。VITS把耦合层限定为只平移、不缩放的保体积变换，行列式恒为1，上式实际只剩高斯密度一项。只用对角高斯作先验时，同一音素各帧的分布完全相同，表达不出帧间变化；流让先验可以是复杂分布，缩小了与后验的差距。第四部分是随机时长预测器，下面单独说明。

对齐由\term{单调对齐搜索}（monotonic alignment search, MAS）求得，该方法最早用于Glow-TTS\cite{nlp-speech-kim2020glowtts}。对齐$A$把每帧$j$分配给一个音素$A(j)$，要求单调且不跳过：第1帧属于第1个音素，最后一帧属于最后一个音素，每帧要么留在前一帧的音素，要么前进一个。MAS在所有合法对齐中选使变换后隐变量在先验下似然最大的一个。记$\ell_{i,j}=\log\mathcal{N}(f_\theta(\bm{z}_j);\bm{\mu}_i,\bm{\sigma}_i)$，用动态规划
\begin{equation}
  Q_{i,j}=\max\bigl(Q_{i-1,j-1},\,Q_{i,j-1}\bigr)+\ell_{i,j},
  \qquad Q_{1,1}=\ell_{1,1},
  \label{eq:nlp-speech-mas}
\end{equation}
填表，再从$Q_{N,T}$回溯得到最优路径，复杂度$O(NT)$。它和第\ref{sec:nlp-classical-viterbi}小节的Viterbi算法同构，只是转移只允许“停留”和“前进一格”。教学示例取2个音素、3帧，$\ell_{1,\cdot}=(-1,-2,-4)$，$\ell_{2,\cdot}=(-5,-1,-1)$。第一行$Q_{1,\cdot}=(-1,-3,-7)$；第二行$Q_{2,1}=-\infty$（第1帧必须属于第1个音素），$Q_{2,2}=\max(-1,-\infty)-1=-2$，$Q_{2,3}=\max(Q_{1,2},Q_{2,2})-1=\max(-3,-2)-1=-3$。回溯得第1帧属于音素1、第2、3帧属于音素2，时长为$(1,2)$，总分$-3$；另一条合法路径（时长$(2,1)$）的总分为$-1-2-1=-4$。训练时每步先固定参数用MAS求对齐，再在该对齐下更新参数，与硬EM（Viterbi训练）的交替方式类似；对齐不参与梯度。

MAS给出的时长$d_i=\sum_j\mathbb{1}[A(j)=i]$还要用来训练一个推断时使用的时长预测器。FastSpeech 2的确定性预测器对同一文字总给出同一节奏。VITS的\term{随机时长预测器}（stochastic duration predictor）是一个基于流的条件生成模型，学习时长的分布$p(d\mid c)$。整数时长不能直接用连续密度建模，模型加入$[0,1)$内的均匀扰动$u$做变分去量化，并引入辅助变量$\nu$增加表达力，最大化
\begin{equation}
  \log p_\theta(d\mid c)\geq\mathbb{E}_{q_\phi(u,\nu\mid d,c)}\Bigl[\log\frac{p_\theta(d-u,\nu\mid c)}{q_\phi(u,\nu\mid d,c)}\Bigr],
  \label{eq:nlp-speech-vits-duration}
\end{equation}
其负值记为$L_{\mathrm{dur}}$。推断时从高斯噪声出发，经逆向流得到连续时长，再向上取整。同一文字采样不同噪声，就得到不同的节奏。

总目标把VAE的两项与对抗训练相加：
\begin{equation}
\begin{aligned}
  L_{\mathrm{VITS}}={}&\lambda_{\mathrm{mel}}\lVert\phi(x)-\phi(\hat x)\rVert_1
  +\underbrace{\mathbb{E}_{q_\phi}\bigl[\log q_\phi(\bm{z}\mid\bm{x}_{\mathrm{lin}})-\log p_\theta(\bm{z}\mid c,A)\bigr]}_{L_{\mathrm{kl}}}\\
  &+L_{\mathrm{dur}}+L_{\mathrm{adv}}(G)+L_{\mathrm{fm}}(G),
\end{aligned}
  \label{eq:nlp-speech-vits-loss}
\end{equation}
第一项是重建项，用梅尔$L_1$误差代替波形似然，相当于假设数据服从拉普拉斯分布；原文把各项直接相加，公开实现与HiFi-GAN一样取$\lambda_{\mathrm{mel}}=45$。第二项是单样本估计的KL项，连接后验与文本先验。对抗与特征匹配项沿用式\eqref{eq:nlp-speech-hifigan}；为节省训练开销，判别器保留HiFi-GAN的5个多周期子判别器，只从多尺度判别器中保留作用于原始波形的一个，相当于周期取$\{1,2,3,5,7,11\}$的多周期判别器。推断时不再需要后验编码器：文本编码器给出$\bm{\mu},\bm{\sigma}$，按随机时长展开到帧，从先验采样$f_\theta(\bm{z})$，经逆向流得到$\bm{z}$，再由解码器生成波形。

VITS的优点是单阶段、并行，训练时就让声码器部分见到模型自己的隐变量，消除了两阶段的不匹配，并且不需要外部对齐器，音质在单说话者数据上接近真人录音。缺点是模型和目标复杂，训练成本高；韵律主要由随机时长和隐变量采样决定，不如FastSpeech 2的显式音高、能量条件便于直接调节。

\subsection{基于扩散与流匹配的合成}
\label{sec:nlp-speech-diffusion-tts}

回归式声学模型的根本问题是一对多。设某句话的句末音高可以上扬（取值$+1$）或下降（取值$-1$），两种读法各占一半；以均方误差训练时，最优预测是条件均值0，一条两种读法都不是的平直曲线。FastSpeech 2的音高预测器、Tacotron 2的频谱输出都受此影响，声音因此显得平淡。对抗训练能避免平均，但不稳定且难以覆盖全部模式。扩散与流匹配模型直接学习条件分布本身：从高斯噪声出发，按学得的方向场逐步把样本移到数据分布上，推断时不同的初始噪声给出不同但都合理的读法。其一般推导见第\ref{sec:nlp-speech-vits}小节所引的卷六章，本小节只写条件构造、目标函数和采样步数的权衡。

\paragraph{Grad-TTS：以文本为均值的分数扩散}
Grad-TTS用扩散模型生成梅尔频谱\scite{nlp-speech-popov2021gradtts}。它的编码器部分与Glow-TTS相同：文本编码器为每个音素输出均值$\tilde{\bm{\mu}}_i\in\mathbb{R}^{80}$，训练时用MAS对齐，再展开为帧级均值$\bm{\mu}\in\mathbb{R}^{T\times80}$；推断时用时长预测器展开。与常见扩散模型把数据扩散到标准高斯不同，Grad-TTS的前向过程把频谱$\bm{X}_0$扩散到以文本均值为中心的高斯$\mathcal{N}(\bm{\mu},\bm{I})$：
\begin{equation}
  \mathrm{d}\bm{X}_t=\tfrac12(\bm{\mu}-\bm{X}_t)\,\beta_t\,\mathrm{d}t+\sqrt{\beta_t}\,\mathrm{d}\bm{W}_t,
  \qquad t\in[0,1],
  \label{eq:nlp-speech-gradtts-forward}
\end{equation}
$\beta_t$是噪声调度，$\bm{W}_t$是布朗运动。这个线性随机微分方程有闭式解：给定$\bm{X}_0$，$\bm{X}_t$服从高斯分布，均值为$\bm{\rho}_t=(1-e^{-\frac12\int_0^t\beta_s\mathrm{d}s})\bm{\mu}+e^{-\frac12\int_0^t\beta_s\mathrm{d}s}\bm{X}_0$，协方差为$\lambda_t\bm{I}$，$\lambda_t=1-e^{-\int_0^t\beta_s\mathrm{d}s}$。于是可以一步采样$\bm{X}_t=\bm{\rho}_t+\sqrt{\lambda_t}\bm{\xi}$，$\bm{\xi}\sim\mathcal{N}(\bm{0},\bm{I})$。网络$\bm{s}_\theta(\bm{X}_t,\bm{\mu},t)$（一个U-Net）学习对数密度的梯度，即分数函数；按去噪分数匹配，目标为
\begin{equation}
  L_{\mathrm{diff}}=\mathbb{E}_{\bm{X}_0,t,\bm{\xi}}\Bigl[\lambda_t\Bigl\lVert\bm{s}_\theta(\bm{X}_t,\bm{\mu},t)+\frac{\bm{\xi}}{\sqrt{\lambda_t}}\Bigr\rVert_2^2\Bigr],
  \label{eq:nlp-speech-gradtts-loss}
\end{equation}
$-\bm{\xi}/\sqrt{\lambda_t}$正是高斯条件密度在$\bm{X}_t$处的分数。总目标再加上编码器输出与频谱之间的高斯负对数似然（让$\bm{\mu}$本身就接近频谱）和时长预测损失。推断从$\mathcal{N}(\bm{\mu},\tau^{-1}\bm{I})$采样$\bm{X}_1$（$\tau$为温度），沿反向常微分方程
\[
  \mathrm{d}\bm{X}_t=\tfrac12\bigl(\bm{\mu}-\bm{X}_t-\bm{s}_\theta(\bm{X}_t,\bm{\mu},t)\bigr)\beta_t\,\mathrm{d}t
\]
从$t=1$积分到$t=0$。起点已经是接近目标的文本均值，而不是纯噪声，因此少量步数就能得到可用的频谱。步数$S$是质量与延迟的权衡：每步要完整运行一次U-Net，延迟随$S$线性增长。

\paragraph{F5-TTS：流匹配与填充式训练}
扩散模型的反向轨迹是弯曲的，需要较多步数才能积分准确。\term{条件流匹配}（conditional flow matching）改为直接回归一个把噪声沿直线运到数据的速度场\cite{lipman2023flowmatching}。取噪声$\bm{x}_0\sim\mathcal{N}(\bm{0},\bm{I})$与数据$\bm{x}_1$，定义直线路径与目标速度
\begin{equation}
\begin{aligned}
  \bm{x}_t&=\bigl(1-(1-\sigma_{\min})t\bigr)\bm{x}_0+t\,\bm{x}_1,\qquad
  \bm{u}_t=\bm{x}_1-(1-\sigma_{\min})\bm{x}_0,\\
  L_{\mathrm{CFM}}&=\mathbb{E}_{t,\bm{x}_0,\bm{x}_1}\bigl\lVert\bm{v}_\theta(\bm{x}_t,t,\text{条件})-\bm{u}_t\bigr\rVert_2^2,
\end{aligned}
  \label{eq:nlp-speech-cfm}
\end{equation}
$\sigma_{\min}$是很小的常数。一维教学示例：取$\sigma_{\min}=0$，$x_0=-0.6$，$x_1=1.0$，$t=0.25$时$x_t=0.75\times(-0.6)+0.25\times1.0=-0.2$，目标速度$u_t=1.6$在整条路径上不变；若网络恰好学到它，从$x_0$出发用欧拉法走两步（步长0.5）得$0.2$、$1.0$，正好到达数据点。真实数据的条件速度场是许多直线的平均，不再处处为常数，但仍比扩散轨迹平直，所需步数更少。

Voicebox把流匹配用于语音的\term{填充式训练}（infilling）：给定一段语音及其文本，遮住其中一部分频谱，让模型根据上下文和文本补全被遮的部分\cite{nlp-speech-le2023voicebox}。合成于是变成填空：把参考录音放在前面、目标部分整个遮住，模型续写出的声音自然沿用参考的音色与风格。F5-TTS进一步去掉了显式时长模型与音素对齐\scite{nlp-speech-chen2024f5tts}。它的输入构造如下：参考文本与目标文本的字符序列拼接后，用填充符补齐到与频谱相同的帧数$T$，经几层ConvNeXt V2块\cite{woo2023convnextv2}得到文本表示$\bm{C}\in\mathbb{R}^{T\times d}$；遮蔽标记$\bm{m}\in\{0,1\}^{T}$标出待生成的帧（训练时随机遮住70\%到100\%的帧），可见部分$(1-\bm{m})\odot\bm{x}_1$作为声音条件。三者与当前噪声频谱$\bm{x}_t$在特征维拼接，送入以时间$t$调制层归一化的扩散Transformer\cite{peebles2023dit}，输出速度$\bm{v}_\theta$。损失只在被遮蔽的帧上计算：
\begin{equation}
  L_{\mathrm{F5}}=\mathbb{E}\bigl\lVert\bm{m}\odot\bigl(\bm{v}_\theta(\bm{x}_t,t,(1-\bm{m})\odot\bm{x}_1,\bm{C})-\bm{u}_t\bigr)\bigr\rVert_2^2 .
  \label{eq:nlp-speech-f5-loss}
\end{equation}
字符与帧之间没有任何对齐监督，模型要在训练中自己学会哪些帧读哪个字；填充符只是让两条序列等长。推断时给出参考录音、参考文本和目标文本，总帧数按目标文本与参考文本的字符数之比由参考时长推算，再从噪声出发用欧拉法积分，原文主要报告32步；还可用无分类器引导，把有条件与无条件的速度外推$\bm{v}=\bm{v}_{\mathrm{cond}}+w(\bm{v}_{\mathrm{cond}}-\bm{v}_{\mathrm{uncond}})$，以加强对文本的服从。生成的频谱交给声码器。

扩散与流匹配模型解决了平均化问题，生成的韵律更富变化；填充式训练使零样本音色迁移只需一段参考录音。代价是多步采样增加了延迟，步数与质量需要折中；F5-TTS一类无显式对齐的模型在长句或罕见文本上可能漏读、重复，总时长靠字符数推算也会让语速偏离参考。

\subsection{基于编解码词元的语言模型合成}
\label{sec:nlp-speech-codec-lm-tts}

扩散与流匹配仍在连续频谱上建模。第\ref{sec:nlp-speech-codec}小节的神经编解码器把波形编码为若干并行的离散码流：残差向量量化（式\eqref{eq:nlp-speech-rvq}）的第一级码本刻画主要内容与音色，后续各级逐级补充残差细节。有了离散表示，合成可以改写为词元预测，直接沿用第\ref{chap:nlp-large-language-models}章大语言模型的训练与采样方法；参考录音的词元放在前面，就成了续写的提示。

VALL-E是这一思路的代表\scite{nlp-speech-wang2023valle}。它使用24\,kHz的EnCodec，每秒75帧，每帧8个码本、每个码本1024个码字，记一段语音的码矩阵为$\bm{C}\in\{1,\ldots,1024\}^{T\times8}$，$c_{t,j}$是第$t$帧第$j$级的码字。1秒声音有$75\times8=600$个词元，3秒提示即1800个。若把8级展平成一条序列自回归生成，序列太长；VALL-E按码本层级拆成两个模型。设目标文本的音素为$x$，提示录音的码矩阵为$\tilde{\bm{C}}$，

\begin{equation}
\begin{aligned}
  p(\bm{C}\mid x,\tilde{\bm{C}})
  ={}&\underbrace{\prod_{t=1}^{T}p_{\mathrm{AR}}\bigl(c_{t,1}\mid c_{<t,1},\,\tilde{\bm{C}}_{:,1},\,x\bigr)}_{\text{第一级：自回归}}
  \;\underbrace{\prod_{j=2}^{8}p_{\mathrm{NAR}}\bigl(\bm{C}_{:,j}\mid\bm{C}_{:,<j},\,\tilde{\bm{C}},\,x\bigr)}_{\text{其余各级：非自回归}} .
\end{aligned}
\label{eq:nlp-speech-valle}
\end{equation}
自回归模型是只含解码器的Transformer，输入序列为“音素$x$，提示的第一级码字$\tilde{\bm{C}}_{:,1}$，已生成的第一级码字”，逐帧预测$c_{t,1}$，直到输出结束符；生成多少帧由模型自己决定，因此不需要时长模型。非自回归模型对第$j$级一次性并行预测所有帧：每帧的输入是前$j-1$级码字嵌入之和，加上音素与完整提示，当前层级$j$通过自适应层归一化注入；8级只需额外7次前向。两个模型都以各级词元的交叉熵训练，与语言模型的下一词元预测相同。训练数据是约6万小时、以识别器自动转写的英语有声书录音。

推断时提供约3秒的参考录音及其转写，把参考转写与目标文本的音素拼接为$x$，参考录音的码矩阵作为$\tilde{\bm{C}}$。自回归模型从提示之后续写第一级码字，采用随机采样；束搜索容易陷入无限重复，原文因此改用采样。非自回归模型补全其余7级，最后由EnCodec解码器把码矩阵变回波形。续写出的声音沿用提示中的说话者、录音环境甚至情绪，这就是\term{零样本语音合成}（zero-shot text-to-speech）：目标说话者从未出现在训练集中，只靠上下文提示完成音色迁移。

与两阶段系统相比，编解码词元语言模型有三点不同。第一，它把音色控制变成上下文续写，不必为每个说话者训练或微调模型，也不必单独设计说话者嵌入。第二，它能沿用语言模型的规模化经验，数据从几十小时的精标录音扩大到数万小时的自动转写语音。第三，目标是离散分布上的交叉熵，采样天然给出多样的韵律，不存在回归平均化。代价同样来自自回归与隐式对齐：采样可能跳词、重复或在句末不停止，内容错误比FastSpeech 2一类显式时长方法更难控制；生成速度受第一级逐帧串行的限制；与参考录音高度相似的声音也带来冒用身份的风险，需要配合水印或说话者授权检查。此后的方案常把两类方法组合：CosyVoice先用语言模型生成从识别监督学到的语义词元，再由流匹配模型把语义词元转为频谱\scite{nlp-speech-du2024cosyvoice}，语言模型负责内容与韵律的离散规划，流匹配负责连续的声学细节。离散声学词元同样是第\ref{sec:nlp-speech-realtime-dialogue}节语音大模型的输出接口。

\begin{table*}[htbp]
  \centering
  \begin{tabularx}{\linewidth}{@{}>{\raggedright\arraybackslash}p{22mm}>{\raggedright\arraybackslash}p{26mm}>{\raggedright\arraybackslash}p{17mm}>{\raggedright\arraybackslash}X>{\raggedright\arraybackslash}p{18mm}>{\raggedright\arraybackslash}X@{}}
    \toprule
    方法 & 对齐方式 & 自回归 & 并行度与推断代价 & 零样本 & 内容错误风险\\
    \midrule
    Tacotron 2 & 位置敏感注意力，隐式 & 逐帧 & 帧数次串行；另需声码器 & 否 & 注意力跳读、重复、不停止\\
    FastSpeech 2 & 外部强制对齐的显式时长 & 否 & 一次前向并行；另需声码器 & 否 & 低；依赖对齐器质量\\
    VITS & 训练中单调对齐搜索 & 否 & 一次前向直接得到波形 & 否 & 低；单调约束保证不跳读\\
    Grad-TTS & 单调对齐搜索与时长 & 否 & 多步去噪，延迟随步数增长 & 否 & 低；对齐显式\\
    F5-TTS & 无显式对齐，填充符等长 & 否 & 约数十步流匹配积分 & 是，参考录音填充 & 中；长句可能漏读\\
    VALL-E & 无显式对齐，模型决定长度 & 首层逐帧 & 首层帧数次串行，其余7级并行 & 是，3秒提示 & 较高；跳词、重复、不停止\\
    \bottomrule
  \end{tabularx}
  \caption{语音合成方法对照。自上而下，对齐由隐式注意力走向显式时长，再由单调搜索走回不显式对齐的大规模生成；零样本能力来自填充式训练或上下文提示。内容错误风险一栏是按机制作出的定性判断，实际程度依赖数据规模与解码设置，需在同一测试文本上用内容核验比较。}
  \label{tab:nlp-speech-tts-methods}
\end{table*}
\FloatBarrier

不论采用哪一类方法，交付的声音都要按前文的四个字段核验内容。若只把已合成通知中的“十七日”改成“十八日”，还需声明修改区间与需要保持的前后区间：新日期时长不同，后续边界要重新对齐，不能把局部替换后的旧时间戳原样保留。音色相似、频谱接近或自然度高，都不能替代内容核验。
